%! Mean, Variance, and Covariance — Unit 8, Lesson 8.4 — Warm-Up
\documentclass[10pt]{article}
\usepackage{linalg-article}
\usepackage{linalg-boxes}
\newcommand{\TallMath}[1]{$\displaystyle #1\rule[-1.4em]{0pt}{3.2em}$}

\begin{document}

\pageheader{Unit 8, Lesson 8.4}{Warm-Up}
\namedateperiod

\vspace{0.12in}

\begin{notesbox}{Getting Ready: Center, Spread, and a Symmetric Matrix}
\small Today we measure the \emph{shape} of data --- where it centers, how it spreads, and how features move
together. That rests on three familiar moves: an average, a sum of squared distances, and the eigenvalues of a
symmetric matrix. Warm up on each.

\vspace{4pt}
\textbf{1. An average (Unit 1).} Four students score $2,4,6,8$ on a quiz. Their \emph{mean} score is
\[
  \mu=\frac{2+4+6+8}{4}=\blank{1.4cm}.
\]
The mean is the center (balance point) of the data.

\vspace{4pt}
\textbf{2. A sum of squared distances (Unit 4).} Measure how far each score sits from the mean $\mu=5$, square
each distance, and add:
\[
  (2-5)^2+(4-5)^2+(6-5)^2+(8-5)^2=\blank{1.6cm}.
\]
Dividing this by $N=4$ will be our measure of \emph{spread}.

\vspace{4pt}
\textbf{3. Eigenvalues of a symmetric matrix (Unit 6).} Find the eigenvalues of
$V=\left[\begin{smallmatrix}5&3\\3&5\end{smallmatrix}\right]$ by solving $\det(V-\lambda I)=0$:
\[
  (5-\lambda)^2-9=0 \;\Rightarrow\; \lambda=\blank{1.0cm}\ \text{or}\ \blank{1.0cm}.
\]
\end{notesbox}

\end{document}
